\documentclass[11pt]{article}
\usepackage[margin=1in]{geometry}
\usepackage{amsmath,amssymb}
\usepackage{graphicx}
\usepackage[colorlinks=true,linkcolor=blue,urlcolor=blue,citecolor=blue]{hyperref}
\setlength{\parindent}{0pt}
\setlength{\parskip}{0.7em}

\title{The missing Poletsky--Stessin bidisk cyclicity range\\
\large An exact later-literature answer}
\author{Run \texttt{fa\_banach\_001} --- \texttt{agent\_lane\_09}}
\date{9 August 2026}

\begin{document}
\maketitle

\textbf{Status.} \texttt{literature\_already\_answered}.  This is a
literature-identification note, not a claim of a new proof.

\section*{Source question}

Torkinejad Ziarati's \emph{Cyclicity in Poletsky--Stessin Weighted Bergman
Spaces}, arXiv:2507.03767, Open Problem 1 on page 18, asks:

\begin{quote}
For $\frac32<\beta\leq2$, determine whether the polynomial
\[
 f(z_1,z_2)=1-\frac{z_1+z_2}{2}
\]
is cyclic in $\mathcal D_\beta(\mathbb D^2)$.
\end{quote}

Here the relevant norm is equivalent to
\[
 \left\|\sum_{k,l\geq0}a_{k,l}z_1^kz_2^l\right\|_\beta^2
 =\sum_{k,l\geq0}(k+l+2)^\beta |a_{k,l}|^2.
\]

\begin{center}
\includegraphics[width=\textwidth]{figures/open_problem_crop.png}
\end{center}

\section*{Exact supporting answer}

Nailwal and Zalar, \emph{Cyclic polynomials in Dirichlet-type Spaces of the
unit bidisk}, arXiv:2511.13441, explicitly state in their abstract that they
solve this problem.  Their Question 1.2 restates it, and Theorem 3.1 proves:
\begin{quote}
If $\alpha\leq2$, then $2-z_1-z_2$ is cyclic in the sum-weight
Dirichlet-type space $\mathcal D_\alpha$.
\end{quote}

The supporting paper uses weights $(k+l+1)^\alpha$, while the source uses
$(k+l+2)^\alpha$; the paper explicitly notes that these norms are equivalent.
Moreover,
\[
 1-\frac{z_1+z_2}{2}=\frac12(2-z_1-z_2),
\]
and multiplication by a nonzero scalar does not change the closed
shift-invariant subspace generated by a polynomial.  Thus Theorem 3.1 gives
a full affirmative answer to the exact source problem, including the endpoint
$\beta=2$.

\begin{center}
\includegraphics[width=\textwidth]{figures/supporting_answer_crop.png}
\end{center}

Their Theorem 1.3 is stronger: for $1<\alpha\leq2$, an irreducible polynomial
without zeros in the bidisk is cyclic precisely when its zero set on
$\mathbb T^2$ is empty or finite.

\section*{Search and scope}

The local run indexes and bounded web searches used arXiv:2507.03767, the
exact polynomial, the paper title, \texttt{Poletsky--Stessin}, and the missing
parameter range.  The later answer was found through a 2026
cyclic-polynomial reference and verified directly in arXiv:2511.13441.  The
supporting authors explicitly know they are solving the source problem.

This identification answers Open Problem 1 only.  The source's Open Problem 2
about one-dimensional boundary zero sets on complex ellipsoids is not covered.

\section*{References}

\begin{enumerate}
\item P. Torkinejad Ziarati, \emph{Cyclicity in Poletsky--Stessin Weighted
Bergman Spaces},\\ arXiv:2507.03767 (2025).
\item R. Nailwal and A. Zalar, \emph{Cyclic polynomials in Dirichlet-type
Spaces of the unit bidisk},\\ arXiv:2511.13441 (2025; revised 2026).
\end{enumerate}

\end{document}
